% GENERATED FILE. Edit canonical lessons, then run npm run generate:books.
% canonical-source-hash: fnv1a64:bf4a94e3b57467c7
% canonical-lessons: RU-C530-stirka, RU-C530-prishchepka, RU-C530-pyatno, RU-C530-razveshivat, RU-C530-otstiryvat

\chapter{Washing Day}
\label{ch:ru-washing-day}
\hlchaptermodality{530}

\begin{chapteropening}
\textbf{By the end of this chapter:} \emph{I can talk about doing the laundry and getting stains out.}

Complete the last lesson of chapter 530: I can talk about doing the laundry and getting stains out.
\end{chapteropening}

\section[\texorpdfstring{\ru{стирка} (stírka)}{стирка (stírka)}]{\ru{стирка} (stírka) — laundry, washing}

\label{lesson:RU-C530-stirka}



\begin{quote}
Before the new one: say the Russian for a cloth, then the Russian for to wipe.
\end{quote}

\subsection*{You'll want to know: \ru{стирка}}
\textbf{\ru{стирка}} (\emph{stírka}) — \textquotedblleft{}laundry, washing\textquotedblright{}.

\begin{quote}
\ru{Стирка} \ru{занимает} \ru{у} \ru{меня} \ru{целый} \ru{вечер}.
\end{quote}

\begin{quote}
\emph{Stírka zanimáyet u menyá tsélyy vécher.}
\end{quote}

\begin{quote}
Doing the laundry takes me a whole evening.
\end{quote}

Feminine.

\begin{cousinweb}
\emph{\ru{стирка}} is built on \emph{\ru{стирать}} 'to wash clothes', first 'to rub off', from the root of \emph{\ru{тереть}} 'to rub'; EtymDB has no entry.
\end{cousinweb}

\begin{grammarlens}[title={What you've built}]
One more word for telling, explaining and describing at B1.
\end{grammarlens}

\subsection*{Your turn}
\begin{itemize}
\raggedright
  \item \emph{Say it:} \emph{stírka}
  \item \emph{Say it:} \emph{stírka}, once more
  \item \emph{Say it:} a sentence of your own with \emph{stírka}
  \item \emph{Recall:} say the Russian for a cloth, then the Russian for to wipe, then say \emph{stírka} again
\end{itemize}

\begin{tcolorbox}[breakable,colback=teal!4,colframe=teal!35!black,title={Before you move on}]
What does \ru{стирка} mean? (\textquotedblleft{}Laundry, washing\textquotedblright{}.) Say it once more.
\end{tcolorbox}
\section[\texorpdfstring{\ru{прищепка} (prishchépka)}{прищепка (prishchépka)}]{\ru{прищепка} (prishchépka) — a clothes peg}

\label{lesson:RU-C530-prishchepka}



\begin{quote}
Before the new one: say the Russian for to wipe, then the Russian for laundry.
\end{quote}

\subsection*{You'll want to know: \ru{прищепка}}
\textbf{\ru{прищепка}} (\emph{prishchépka}) — \textquotedblleft{}a clothes peg\textquotedblright{}.

\begin{quote}
\ru{На} \ru{верёвке} \ru{висит} \ru{бельё} \ru{на} \ru{прищепках}.
\end{quote}

\begin{quote}
\emph{Na veryóvke visít belyó na prishchépkakh.}
\end{quote}

\begin{quote}
The washing is hanging on the line with pegs.
\end{quote}

Feminine.

\begin{cousinweb}
\emph{\ru{прищепка}} is built on \emph{\ru{прищепить}} 'to pin on, to clip', usually linked to \emph{\ru{щепа}} 'a chip of wood'; EtymDB has no entry.
\end{cousinweb}

\begin{grammarlens}[title={What you've built}]
One more word for telling, explaining and describing at B1.
\end{grammarlens}

\subsection*{Your turn}
\begin{itemize}
\raggedright
  \item \emph{Say it:} \emph{prishchépka}
  \item \emph{Say it:} \emph{prishchépka}, once more
  \item \emph{Say it:} a sentence of your own with \emph{prishchépka}
  \item \emph{Recall:} say the Russian for to wipe, then the Russian for laundry, then say \emph{prishchépka} again
\end{itemize}

\begin{tcolorbox}[breakable,colback=teal!4,colframe=teal!35!black,title={Before you move on}]
What does \ru{прищепка} mean? (\textquotedblleft{}A clothes peg\textquotedblright{}.) Say it once more.
\end{tcolorbox}
\section[\texorpdfstring{\ru{пятно} (pyatnó)}{пятно (pyatnó)}]{\ru{пятно} (pyatnó) — a stain, a spot}

\label{lesson:RU-C530-pyatno}



\begin{quote}
Before the new one: say the Russian for laundry, then the Russian for a clothes peg.
\end{quote}

\subsection*{You'll want to know: \ru{пятно}}
\textbf{\ru{пятно}} (\emph{pyatnó}) — \textquotedblleft{}a stain, a spot\textquotedblright{}.

\begin{quote}
\ru{На} \ru{рубашке} \ru{осталось} \ru{пятно} \ru{от} \ru{кофе}.
\end{quote}

\begin{quote}
\emph{Na rubáshke ostálos pyatnó ot kófe.}
\end{quote}

\begin{quote}
There's a coffee stain left on the shirt.
\end{quote}

Neuter; plural \ru{пятна}.

\begin{cousinweb}
EtymDB lists \emph{\ru{пятно}} without a source, and the origin is not settled.
\end{cousinweb}

\begin{grammarlens}[title={What you've built}]
One more word for telling, explaining and describing at B1.
\end{grammarlens}

\subsection*{Your turn}
\begin{itemize}
\raggedright
  \item \emph{Say it:} \emph{pyatnó}
  \item \emph{Say it:} \emph{pyatnó}, once more
  \item \emph{Say it:} a sentence of your own with \emph{pyatnó}
  \item \emph{Recall:} say the Russian for laundry, then the Russian for a clothes peg, then say \emph{pyatnó} again
\end{itemize}

\begin{tcolorbox}[breakable,colback=teal!4,colframe=teal!35!black,title={Before you move on}]
What does \ru{пятно} mean? (\textquotedblleft{}A stain, a spot\textquotedblright{}.) Say it once more.
\end{tcolorbox}
\section[\texorpdfstring{\ru{развешивать} (razvéshivat)}{развешивать (razvéshivat)}]{\ru{развешивать} (razvéshivat) — to hang out (washing)}

\label{lesson:RU-C530-razveshivat}



\begin{quote}
Before the new one: say the Russian for a clothes peg, then the Russian for a stain.
\end{quote}

\subsection*{You'll want to know: \ru{развешивать}}
\textbf{\ru{развешивать}} (\emph{razvéshivat}) — \textquotedblleft{}to hang out (washing)\textquotedblright{}.

\begin{quote}
\ru{Мама} \ru{развешивает} \ru{бельё} \ru{на} \ru{балконе}.
\end{quote}

\begin{quote}
\emph{Máma razvéshivayet belyó na balkóne.}
\end{quote}

\begin{quote}
Mum is hanging out the washing on the balcony.
\end{quote}

Imperfective; perfective \ru{развесить} (\ru{развешу}).

\begin{cousinweb}
\emph{\ru{развешивать}} is built on \emph{\ru{раз}-} and \emph{\ru{вешать}} 'to hang'; EtymDB has no entry.
\end{cousinweb}

\begin{grammarlens}[title={What you've built}]
One more word for telling, explaining and describing at B1.
\end{grammarlens}

\subsection*{Your turn}
\begin{itemize}
\raggedright
  \item \emph{Say it:} \emph{razvéshivat}
  \item \emph{Say it:} \emph{razvéshivat}, once more
  \item \emph{Say it:} a sentence of your own with \emph{razvéshivat}
  \item \emph{Recall:} say the Russian for a clothes peg, then the Russian for a stain, then say \emph{razvéshivat} again
\end{itemize}

\begin{tcolorbox}[breakable,colback=teal!4,colframe=teal!35!black,title={Before you move on}]
What does \ru{развешивать} mean? (\textquotedblleft{}To hang out (washing)\textquotedblright{}.) Say it once more.
\end{tcolorbox}
\section[\texorpdfstring{\ru{отстирывать} (otstíryvat)}{отстирывать (otstíryvat)}]{\ru{отстирывать} (otstíryvat) — to wash out (a stain)}

\label{lesson:RU-C530-otstiryvat}



\begin{quote}
Before the new one: say the Russian for a stain, then the Russian for to hang out.
\end{quote}

\subsection*{You'll want to know: \ru{отстирывать}}
\textbf{\ru{отстирывать}} (\emph{otstíryvat}) — \textquotedblleft{}to wash out (a stain)\textquotedblright{}.

\begin{quote}
\ru{Пятно} \ru{от} \ru{вина} \ru{трудно} \ru{отстирывать}.
\end{quote}

\begin{quote}
\emph{Pyatnó ot viná trúdno otstíryvat.}
\end{quote}

\begin{quote}
A wine stain is hard to wash out.
\end{quote}

Imperfective; perfective \ru{отстирать}.

\begin{cousinweb}
\emph{\ru{отстирывать}} is built on \emph{\ru{от}-} 'off' and \emph{\ru{стирать}} 'to wash clothes', from the root of \emph{\ru{тереть}} 'to rub'; EtymDB has no entry.
\end{cousinweb}

\begin{grammarlens}[title={What you've built}]
One more word for telling, explaining and describing at B1.
\end{grammarlens}

\subsection*{Your turn}
\begin{itemize}
\raggedright
  \item \emph{Say it:} \emph{otstíryvat}
  \item \emph{Say it:} \emph{otstíryvat}, once more
  \item \emph{Say it:} a sentence of your own with \emph{otstíryvat}
  \item \emph{Recall:} say the Russian for a stain, then the Russian for to hang out, then say \emph{otstíryvat} again
\end{itemize}

\begin{tcolorbox}[breakable,colback=teal!4,colframe=teal!35!black,title={Before you move on}]
What does \ru{отстирывать} mean? (\textquotedblleft{}To wash out (a stain)\textquotedblright{}.) Say it once more.
\end{tcolorbox}
